% !TeX root = book.tex
The rapid growth of e-commerce and smart logistics has increased the demand for efficient last-mile delivery solutions. Drone delivery has emerged as a promising method for transporting lightweight packages in environments such as smart campuses, industrial parks, residential areas, and remote locations. However, effective drone delivery requires more than drone operation alone. It involves managing customers, delivery requests, packages, landing stations, delivery tracking, and confirmation through a centralized and reliable information system. Existing logistics systems are mainly designed for traditional transportation and often lack dedicated support for drone-based workflows. To address these challenges, the SmartDroneDelivery project proposes an AI-powered Drone Delivery Management Platform that enables efficient management, coordination, and monitoring of drone delivery operations in a business environment.
\section{SmartDroneDelivery Platform }\label{s:introduction} 

SmartDroneDelivery is designed to support the complete delivery lifecycle, from the creation of a delivery request to the confirmation of successful delivery. The platform focuses on business and logistics management rather than direct drone control.

The system supports five primary business processes:

\begin{enumerate}
\item Delivery request management
\item Package management
\item Delivery execution
\item Delivery tracking
\item Delivery confirmation
\end{enumerate}

In addition to these core processes, AI-assisted services are proposed to improve operational efficiency and customer experience. These services include estimated delivery time, delivery summarization, customer assistance, and operational analytics.
The platform supports five main human actors:

\begin{enumerate}
\item Customer
\item Dispatcher
\item Station Operator
\item Logistics Manager
\item System Administrator
\end{enumerate}

The system is organized into five core modules:

\begin{enumerate}
\item Customer and Delivery Order Management
\item Package and Delivery Workflow Management
\item Landing Station Management
\item Delivery Tracking and Confirmation
\item Analytics and AI-assisted Services
\end{enumerate}

The proposed platform provides a web portal for logistics management, a mobile application for customers and operators, and standardized RESTful APIs for integration with external services.
\section{Database Perspective} 

\qquad Within the SmartDroneDelivery platform, the database plays a central role in storing and managing information generated by the different business processes. The database must support the management of users, customers, addresses, delivery orders, packages, drones, landing stations, delivery activities, tracking information, delivery confirmation, and other supporting information required by the system.

The database design must ensure that data is organized in a structured manner and that relationships among different entities are clearly defined. It must also maintain data integrity and consistency through appropriate primary keys, foreign keys, constraints, and other database mechanisms.

A well-designed database provides a reliable data foundation for the backend services and allows the frontend and other system components to access and process data consistently. Therefore, database design is an essential part of the SmartDroneDelivery platform.
\section{Scope of This Document} 

This document focuses specifically on the database design of the SmartDroneDelivery platform.
The main objectives of the database design are:

\begin{enumerate}
\item Analyze the data requirements derived from the proposed business processes.
\item Identify the main entities and their attributes.
\item Define relationships and cardinalities among entities.
\item Develop the conceptual database model.
\item Transform the conceptual model into a logical relational schema.
\item Apply database normalization principles to reduce redundancy and improve consistency.
\item Define primary keys, foreign keys, and integrity constraints.
\item Design indexes and other database structures to support efficient data access.
\item Implement the database using PostgreSQL.
\item Provide a consistent database foundation for future backend and frontend integration.
\end{enumerate}

The database design presented in this document is therefore intended to serve as the data foundation of the SmartDroneDelivery platform and to support the future implementation of its web, mobile, API, and AI-assisted services.
\section{Document Organization} 

\qquad The document is organized into several sections covering the analysis and design of the SmartDroneDelivery database.
The introductory section presents the project background, system overview, database perspective, and scope of the database design.
The following sections describe the data requirements and business rules, followed by the conceptual database model, logical relational schema, normalization, integrity constraints, physical database design, implementation, and testing.
Finally, the document presents the database implementation results and provides conclusions regarding the proposed database design.

